% Unite: papier 34, continuation et section 5; traduction francaise depuis la source allemande auditee.

\providecommand{\mG}{\mathfrak{G}}
\providecommand{\mA}{\mathfrak{A}}
\providecommand{\mE}{\mathfrak{E}}
\providecommand{\mH}{\mathfrak{H}}

\subsection*{\S{} 5. Groupes complètement réductibles}

Un groupe s'appelle \emph{complètement réductible} lorsqu'il est le produit direct d'un nombre fini de groupes simples:
\[
        \mG=\mA_1\times\cdots\times\mA_n.
\]
Dans ce cas, les groupes
\[
        \mA_1\times\cdots\times\mA_n
        \supset \mA_1\times\cdots\times\mA_{n-1}
        \supset\cdots\supset \mA_1\supset \mE
\]
forment une suite de composition. La longueur est $n$, et les facteurs de composition sont
\[
        \sim \mA_n,\ \mA_{n-1},\ldots,\mA_1.
\]

\emph{Si $\mG$ est complètement réductible, alors tout diviseur normal est facteur direct, et l'autre facteur peut être choisi comme produit de groupes simples qui apparaissent dans une décomposition en produit donnée de $\mG$. Le diviseur normal $\mH$ est lui-même complètement réductible.}

\emph{Démonstration.} On a
\begin{equation}
        \mG=\mH\mA_1\mA_2\cdots\mA_n.             \tag{3}
\end{equation}
Si l'on pose maintenant $\mH_i=\mH\mA_1\cdots\mA_i$, alors $\mH_{i+1}=\mH_i\mA_{i+1}$. Ou bien $\mA_{i+1}\leq \mH_i$, donc $\mH_{i+1}=\mH_i$; nous omettons alors dans (3) le facteur $\mA_{i+1}$. Ou bien $\mH_i\cap\mA_{i+1}$ est un diviseur normal propre de $\mA_{i+1}$, donc égal à $\mE$, et alors $\mH_i\times\mA_{i+1}$ est direct. Après omission des facteurs superflus, (3) devient
\[
        \mG=\mH\times\mA_{i_1}\times\cdots\times\mA_{i_r},
\]
donc $\mH$ est facteur direct. En outre, on en déduit que
\[
        \mH\sim \mG/(\mA_{i_1}\times\cdots\times\mA_{i_r})
        \sim \mA_{j_1}\times\cdots\times\mA_{j_\mu},
\]
où $\mA_{j_1},\ldots,\mA_{j_\mu}$ sont les autres $\mA_i$; donc $\mH$ est complètement réductible.

\emph{Corollaire.} \emph{Toute décomposition d'un groupe complètement réductible en facteurs directement indécomposables est une décomposition en facteurs simples, et donne par conséquent lieu à une suite de composition; il en résulte la détermination univoque des facteurs, à isomorphie près.}
